% ============================================================================
% SEC02.TEX - Section 10.2: Membuat Enemy Spawner
% ============================================================================

\section{Membuat Enemy Spawner}

Kita tidak ingin meletakkan asteroid secara manual di Scene (itu membosankan). Kita perlu script yang melahirkannya secara otomatis.

\subsection{Script Spawner}
Buat GameObject kosong bernama \texttt{EnemySpawner} di luar layar (misalnya di atas layar, Y = 6).

\begin{lstlisting}[language={[Sharp]C}]
using UnityEngine;

public class EnemySpawner : MonoBehaviour
{
    public GameObject asteroidPrefab;
    public float spawnRate = 2f; // Detik antar spawn
    private float nextSpawnTime = 0f;
    
    // Area spawn (lebar X)
    public float spawnWidth = 8f;


    void Update()
    {
        if (Time.time >= nextSpawnTime)
        {
            SpawnAsteroid();
            nextSpawnTime = Time.time + spawnRate;
        }
    }

    \begin{figure}[ht]
        \centering
        \begin{tikzpicture}[node distance=1.5cm, auto]
            \node (start) [draw, rounded corners] {Update Loop};
            \node (check) [draw, diamond, aspect=2, below of=start] {Time >= NextSpawn?};
            \node (spawn) [draw, rectangle, below of=check, yshift=-0.5cm] {SpawnAsteroid()};
            \node (reset) [draw, rectangle, below of=spawn] {Reset Timer};
            
            \draw[->] (start) -- (check);
            \draw[->] (check) -- node {Yes} (spawn);
            \draw[->] (check) -- node[right] {No} +(2,0) |- (start);
            \draw[->] (spawn) -- (reset);
            \draw[->] (reset) -- +(2.5,0) |- (start);
        \end{tikzpicture}
        \caption{Logika Spawner}
        \label{fig:spawner_logic}
    \end{figure}

    void SpawnAsteroid()
    {
        // Tentukan posisi X acak
        float randomX = Random.Range(-spawnWidth, spawnWidth);
        Vector3 spawnPos = new Vector3(randomX, transform.position.y, 0);
        
        Instantiate(asteroidPrefab, spawnPos, Quaternion.identity);
    }
}
\end{lstlisting}

\subsection{Setup}
\begin{itemize}
    \item Drag prefab \textbf{Asteroid} ke slot \texttt{Asteroid Prefab}.
    \item Mainkan game. Asteroid akan muncul setiap 2 detik di posisi acak di atas layar.
\end{itemize}
